\documentclass{ltjsbook}
\usepackage{docmute} 
\usepackage{../../myStyle} 

\begin{document}
% \chapter{現代テスト理論その2}
\chapter{Modern Test Theory 2}
\label{x05_GRM}

% \section{現代テスト理論の特徴}
\section{Characteristics of Modern Test Theory}
% 前回は現代テスト理論として項目反応理論(IRT)を紹介しました。ロジスティック曲線を応用して項目の特徴を描画し，それを使って被験者母数を推定する方法についても解説しました。この一連の手続きに基づき，現代テスト理論の利点を考えてみたいと思います。
Last time, I introduced Item Response Theory (IRT) as a modern test theory. I also explained the method of estimating the parameter of the examinee by drawing the characteristics of the item using a logistic curve. Based on this series of procedures, I would like to consider the advantages of modern test theory.
% \subsection{現代テスト理論の利点1:項目母数と被験者母数の分離}
\subsection{Advantage 1 of Modern Test Theory: Separation of Item Parameters and Subject Parameters}
% 古典的テスト理論からの発展として，現代テスト理論では被験者母数$\theta_i$と項目母数$a_j, b_j$を区分して考えるようになりました。項目母数は通過率のアイデアを精緻にしたものですが，この通過率は項目群の総和を元に考えられていたことを思い出してください。
From the development of classical test theory, in modern test theory, we now consider the subject population $\theta_i$ and the item population $a_j, b_j$ separately. Although the item population is a more refined idea of the pass rate, remember that this pass rate was based on the total sum of the item groups.
% すなわち，項目母数の計算には項目の相対的な困難度だけを用いています。イメージとしては鉱物の硬度検査のようなものです。2つの異なる硬さの物をぶつけて崩れた方が負け＝より硬度が低いと考えるように，2つの異なる項目を被験者に与えて，より正答者数が少ない方がより難しいと考えるのです。これはつまり，回答者の学力が高かろうが低かろうが，困難度が$b_x < b_y$であるという関係に違いはないという考え方です。
In other words, the calculation of item parameters only uses the relative difficulty of the items. The image is like a mineral hardness test. As we consider the one that crumbles when hit by two things of different hardnesses as the loser = having lower hardness, when we give two different items to the examinees, we consider the one with fewer correct answers to be more difficult. In other words, this is the idea that regardless of whether the respondent's academic ability is high or low, there is no difference in the relationship that the difficulty is $b_x < b_y$.

% これまでのテストや心理尺度の作り方に比べると，この点が大きく違います。学校などのテストは教員が作っていますから，教員が自分の感覚で「こちらの方がより発展的な内容だ」「こちらの方が難しいだろう」という問いに大きな配点がなされたりするでしょう。その後テストの平均点をみて「今回のテストは簡単にしすぎたか」という判断をしたりするでしょう。しかしIRTでは項目それ自体に困難度を決めさせますから，そこに作成者や回答者の意図は含まれません。平均正答率が高いからと言って簡単な問題なのではなく，項目の特徴として困難度が決まるのです。
This is significantly different from the way previous tests and psychological measures were created. Tests in schools and the like are made by teachers, so they may assign larger scores to questions that, by their own judgment, have more advanced content or seem harder. They may then look at the average test scores and decide whether they made the test too easy. However, with IRT, difficulty is determined by the items themselves, so neither the creator's nor the respondent's intent are included. It's not that a problem is easy because the average correct answer rate is high, rather the difficulty is determined by the characteristics of the item.

% たとえばサーストン尺度の作り方を思い出してください(セクション\ref{ThurstoneScale},Pp.\pageref{ThurstoneScale}参照)。サーストン尺度では尺度適用前に評定者集団によって尺度値を決めます。この評定者集団が偏った思想の持ち主だけで固められていた場合，尺度の点数は極端なものになり，普通の人がその尺度に回答すると極めて低い点数，高い点数になってしまうかもしれません。
For example, remember how to create a Thurston scale (see Section \ref{ThurstoneScale}, Pp.\pageref{ThurstoneScale}). In the Thurston scale, scale values are determined by a panel of judges before the scale is applied. If this group of judges is made up of people with biased thoughts, the scale scores may become extreme and ordinary people answering the scale could score very low or very high.
% あるいはリッカート尺度の作り方を思い出してください。(セクション\ref{LikertScale},Pp.\pageref{LikertScale}参照)。リッカート尺度では回答者の累積度数から尺度値を算出します。先ほど同様，回答者集団の態度に偏りがあれば，尺度の点数は標準的な物ではなくなるでしょう。
Please recall how to make a Likert scale. (Refer to Section \ref{LikertScale}, Pp. \pageref{LikertScale}). In Likert scale, the scale value is calculated from the cumulative frequency of the respondent's answers. Similarly to before, if there is a bias in the attitude of the group of respondents, the score on the scale will no longer be standard.
% つまり「誰を対象に測定するか」によって目盛りが変わるようなものです。これでは測定結果の一般化は難しいでしょう。たとえば本学で作られた尺度を，他大学でやってみると違った尺度値になるのですから，研究結果はせいぜい「その大学ではそうなんだろう」となり一般化できなくなります。
In other words, the scale changes depending on "who is being measured". This makes it difficult to generalize the measurement results. For example, if you try to use the scale created in this university at another university, you will get different scale values, so the research results will be at best "this is probably the case in that university", making it impossible to generalize.
% 従来の方法は，回答者と項目の特徴が関連しすぎていたのです。
The conventional method was too closely related to the characteristics of the responders and the items.

% これに対し，IRTを使って項目の特徴を計算する場合は，相対的な難易度に違いはありませんから，どの大学で作った尺度であっても統一的な解釈が可能です。尺度作成時に幅広くデータを集め，項目母数を確定させてしまえばどこででも統一的な評価ができます。テストなど学力を測定する際に大学間での違いが見られたとしても，その難易度を調整するのも簡単で，共通する項目を入れておけばそこを基準に相対的な困難度調整ができます\footnote{こうしたテスト間の困難度調整のことをテストの\keyterm{等価}{とうか}{equation}といいます。}。良問と悪問の評価と，回答者の評価を分けることは重要なポイントなのです。
In response to this, when using IRT to calculate the characteristics of an item, there is no difference in relative difficulty, so a uniform interpretation is possible regardless of the scale created at any university. If you collect a wide range of data at the time of scale creation and determine the number of items, you can make a uniform evaluation anywhere. Even if differences are observed between universities when measuring academic ability such as tests, it is easy to adjust the difficulty level, and if you include common items, you can adjust the relative difficulty based on that\footnote{This adjustment of difficulty between tests is called the \keyterm{equivalency}{とうか}{equation} of the test.}. It is important to separate the evaluation of good and bad questions and the evaluation of the respondent.
% \subsection{現代テスト理論の利点2:被験者母数の推定の利点}
\subsection{Advantages of Modern Test Theory 2: Advantages of Estimating the Population of Subjects}
% 被験者母数と項目母数の分離は，さらに別の利点も生み出します。それはデータが一部欠落した場合の補完に関係します。
The separation of the subject population and item population produces another advantage. It relates to the completion when some of the data is missing.

% リッカート尺度では回答者全体の相対頻度から，カテゴリの尺度値を決定するのでした。ここである人が特定の項目にだけ回答をし忘れたとします(調査研究ではよくあることで，同じような目盛りが並んでいると一行飛ばして丸をつけちゃうようなことはよくあります)。そうすると，その項目だけ合計人数が変わりますから，計算が面倒です。また相関係数を計算する時にも，その人のその項目については計算できなくなります。因子分析は相関係数から計算を始めますから，一箇所でも欠損値があるとその人のデータを抜いてしまうか\footnote{リストワイズ削除といいます。}，他の値を代入して補完するか\footnote{欠損値補完については，平均値や中央値を代入したり，同じようなパターンで回答している人の値を使い回したり，回帰分析でその項目の値の推定値を入れたり，とさまざまな方法が考えられてきました。欠損値発生のメカニズムにもよりますが，いずれもある程度バイアスのかかった値になってしまいます。統計的によりバイアスの少ない適切な代入法が考えられてはいますが，そもそも欠損値がないのが最も望ましいことに変わりはありません。}，手間でも計算の時にそこだけ外して計算するか\footnote{ペアワイズ削除といいます。}といった工夫が必要です。因子分析結果に大きな違いは出なくても，その人の因子得点は計算できないことに違いはありません。
In a Likert scale, the scale values of categories are determined from the relative frequencies of all respondents. Consider the case where a person forgets to respond to a particular item (which is common in survey research; it often happens that a mark is made skipping a line when similar scales are lined up). Then, the total number of that item changes, making the calculation troublesome. Also, when calculating the correlation coefficient, that person's item cannot be computed. Factor analysis begins with the calculation from the correlation coefficient, so if there is any missing value, either you have to eliminate that person's data \footnote{This is called listwise deletion.}, substitute other values to compensate \footnote{For missing value imputation, various methods have been proposed, such as substituting mean or median values, recycling the values of people who answered in similar patterns, and inserting estimated values for that item from regression analysis. Depending on the mechanism of missing value occurrence, any method results in a value with some bias. Although statistical methods that are less biased have been proposed, it's most desirable to have no missing values in the first place.}, or make an effort to exclude it only during the calculation \footnote{This is called pairwise deletion.}. Although the factor analysis results do not considerably differ, there's no doubt that that person's factor scores cannot be computed.


% それに対して，IRTの被験者母数の推定は，項目母数をつかったICCをもとに一人ずつ絞り込んでいくというものでした。もしある人が回答していないことがあっても，計算ができなくなることはありません。その項目の情報が得られないので絞り込み精度は上がりませんが，ヒントが減っただけで回答できないわけではないのです。このように，得られた情報すべてを使ってその人の被験者母数(因子得点)を推定する方法のことを，\keyterm{完全情報最尤推定}{かんぜんじょうほうさいゆうすいてい}{Full Information Maximum Likelihood}と言います。このようにIRTを使うと，必ずしも全問に回答していなくてもスコアは計算できるということになります。欠損値があるからその人のデータは使えない，ということがないのでいいですね！
In contrast, the estimation of the subject population in IRT was done by narrowing down one by one based on the ICC using the item population. Even if there are some questions that a person did not answer, it does not make the calculation impossible. Because we can't get the information of that item, the precision of narrowing down won't improve, but it doesn't mean you can't answer just because you've lost a clue. The method of estimating a person's subject population (factor score) using all the information obtained in this way is called Full Information Maximum Likelihood. Thus, by using IRT, it is possible to calculate scores even if all questions are not necessarily answered. It's nice that you don't have to say that a person's data can't be used just because there are missing values!

% もっというと，IRTでは全員が全員同じ問題に回答する必要はありません。たとえば能力値が$\theta_i = -0.3$ぐらいにありそうだ，と絞り込めている段階で，次の問題の困難度母数が$b_j = 2.5$であれば，おそらくほぼ確実にその人は回答できないでしょう\footnote{2PLモデルで$a=1,b=2.5$のとき，$\theta=-0.3$が正解する確率は0.8492\%です。}。
In other words, in IRT, it is not necessary for everyone to answer all the questions. For example, if you have narrowed down the ability value to about $\theta_i = -0.3$, and if the difficulty parameter of the next problem is $b_j = 2.5$, that person probably will not be able to answer almost certainly\footnote{In the 2PL model with $a=1, b=2.5$ , the probability that $\theta=-0.3$ will be correct is 0.8492\%.}.
% その人にいくら困難度母数の高い質問を繰り返しても，ほとんど誤答がつづくだけで，とくにその人の$\theta$がありそうな領域を狭めるヒントにはなりません。むしろ$b_j=-1$とか$b_j=-0.5$のような簡単な問題を出して\footnote{2PLモデルで$a=1,b=-1$のとき，$\theta=-0.3$が正解する確率は76.674\%，$b=-0.5$のであれば58.419\%です。}，それらに正答できるかどうかを見極め，絞り込んで行った方が効率的です。
Even if you repeatedly ask this person questions with a high degree of difficulty, they will mostly be answered incorrectly, and this won't particularly provide any hints to narrow down the person's theta. Instead, asking simpler questions like $b_j=-1$ or $b_j=-0.5$ to see if they can answer them correctly and refining based on that would be more efficient. In the 2PL model, when $a=1$, $b=-1$, the probability that $\theta=-0.3$ is correct is 76.674%, and if $b=-0.5$, then it's 58.419%.
% 紙に印刷されているテスト(Paper Based Test)であれば，印刷された問題は変えようがありませんから，困難度順に問題を並べると，あるところから先はずーっと不正解が続く人が続出します。ずーっと不正解なところの問題はいくら良問でも，その人の能力を測るのには役立ちません。ヒントが増えないからです。であれば，回答者の学力に合った問題を，その都度その都度出題した方がいいですよね。コンピュータを使って回答者に相応しい質問をダイナミックに組み替える，コンピュータに基づいたテスト(Computer Based Test)，別名\keyterm{コンピュータ適応型テスト}{こんぴゅーたてきおうがたてすと}{Computer Adaptive Test}というのがそれです。
If it's a paper-based test, once the questions are printed, they cannot be altered, so if the problems are arranged in order of difficulty, beyond a certain point, there would be a succession of people who are continuously incorrect. No matter how good the questions are where they are constantly wrong, they are not useful for measuring the person's abilities. This is because there are no more hints. Therefore, it would be better to provide questions that suit the test-taker's academic abilities each time. A computer-based test, also known as a "Computer Adaptive Test", in which appropriate questions for the respondent are dynamically rearranged using a computer, is just that.
% CATになると，回答者ごとに出題が変わりますからカンニング対策の必要も無くなって，とても便利になること間違いなしです。
When it becomes CAT, the questions change for each respondent, so there's no need for anti-cheating measures, and it's definitely going to be very convenient.

% \subsection{現代テスト理論の利点3:信頼性についての考え方}
\subsection{Advantage 3 of Modern Test Theory: Perspective on Reliability}
% IRTの考え方からいうと，どんな項目でも何らかの情報を提供してくれるはずです。たとえば$b_j=3$のようなとても難しい項目が合ったとします。これがどれぐらい難しいかというと，偏差値70の人でも15\%しか正解できないレベルです。
From the perspective of Item Response Theory (IRT), every item should provide some information. For example, suppose there is a very difficult item, such as $b_j=3$. This difficulty level means that even a person with a deviation value of 70 can only answer correctly 15% of the time.
% ほとんどの人にとっては誤答にしかならない難しすぎる悪問だ，と言いたくなるかもしれませんが，学力が高い人がどれほど高いレベルでやれるのかを検証するためには必要な問題です。偏差値が70なのか，75なのか，80までいけるのか，といったことを見極めるためにはこの問題でないと情報が得られないことになります。逆に簡単すぎる問題であっても，より低いレベルを精緻に検証するためには必要なものなのです。
You might want to say that it's a too difficult question that only leads to wrong answers for most people, but it's a question needed to verify how well people with high academic abilities can perform at a high level. To determine whether the deviation value is 70, 75, or even up to 80, this question is the only way to obtain such information. On the other hand, even if the question is too simple, it is necessary to accurately verify a lower level.

% つまり，どの項目にもその項目が得意とする領域があるはずです。この項目はこの領域の回答者を絞り込む時に，最も有用な情報をもたらしてくれるはず，という$\theta$の場所があるはずなのですね。
In other words, each item should have a domain in which it excels. There should be a place for $\theta$ that brings the most useful information when narrowing down the respondents in this domain.
% これを表現するのが\keyterm{項目情報曲線}{こうもくじょうほうきょくせん}{Item Information Curve}といい，次の式で表される項目情報関数で描くことができます。
This is represented by the \keyterm {Item Information Curve} {こうもくじょうほうきょくせん}, which can be drawn with the item information function represented by the following formula.

\[ I(\theta) = a_j^2 p_j(\theta)q_j(\theta)\]

ここで$p_j(\theta)$はその項目$j$の$\theta$における正答率(通過率)，$q_j(\theta)$は誤答率を表しています。能力が平均的，すなわち$\theta=0$のときは，$0.5 \times 0.5$になる平均的な困難度($b=0$)の設問が最も大きな値になる，というわけですね。
図\ref{fig::05_01}にいくつかのICCとそれに対応するIICを描きましたので，確認してください。
\begin{figure}[ht]
    \begin{center}
        \includegraphics[width=0.8\linewidth,keepaspectratio,clip]{../images/05_GRM/Rplot05_01.png}
        \caption{項目特徴曲線(ICC,上図)と項目情報曲線(IIC,下図)の例。左から順に$a_1=0.8,b_1=-1$の識別力が弱く簡単な項目，$a_2=1,b=1$のやや困難な項目，$a_3=1.2,b_3=2$の識別力が高く困難な項目。}
        \label{fig::05_01}
    \end{center}
\end{figure}
IICのピークは，対応するICCの$\theta=0.5$のところにあること，識別力はピークの尖り具合に関わっていることを確認してください。

さて，IICはその項目がどこで情報をもたらしてくれるか，ということを表しているのでした。言い方を変えると，IICのピークはその項目の最も信頼できるところであるとも言えます。つまりIRTにおいて\textbf{信頼性は潜在特性の関数}になっているのです。
古典的テスト理論ではテスト全体の分散に占める真分散の割合のことを信頼性というのでした。因子分析理論では項目の中の共通因子負荷量の二乗和，共通性がその項目の信頼性を表しているのでした。信頼性を見る水準がテスト全体から項目別の評価に発展したわけですが，IRTではさらにその項目の最も感度の良いところを探る関数として，その信頼性を評価できるようになったといえるでしょう。

またIICはある項目から得られる情報のことを意味しますが，テストに含まれているすべての情報関数を足し合わせることで，そのテストから得られる情報の大きさを関数として評価できます。すなわちテスト全体の情報量$I_T(\theta)$は，$I_T(\theta)=\sum_{j=1}^M I_j(\theta)$であり，この関数のことを\keyterm{テスト情報関数}{てすとじょうほうかんすう}{Test Infromation Curve}といいます。図\ref{fig::05_01}の3つの項目からなるTICを示したのが図\ref{fig::05_02}です。
\begin{figure}[ht]
    \begin{center}
        \includegraphics[width=0.8\linewidth,keepaspectratio,clip]{../images/05_GRM/Rplot05_02.png}
        \caption{テスト情報関数(黒の実線部)。理解を進めるために各IICを薄い点線で表現した。}
        \label{fig::05_02}
    \end{center}
\end{figure}
これを見ると，この3つの項目からなるテストは$\theta=2$より少し低いレベルを測定するときに最も鋭敏に働くということがわかります。このように，項目母数がわかっていれば事前にテストの特徴をどのあたりに持ってくるかを決定でき，自由自在にテストをデザインできるようになるわけです。

テストの例で話をしていますが，心理学的な領域でももちろん便利な手法です。たとえば高い認知能力レベルの人をとくに選出したいとか，精神的な健康度がごく低い人をしっかりと検出したいといった目的があれば，そのあたりにピークが来るような項目からなる質問項目からなる調査票を構成すれば良いのです\footnote{具体例として\citet{kosugi120005596041}をあげておきます。学校適応感を測定するため，テストのピークがやや低いところに来るようになっています。}。


\subsection{現代テスト理論の問題点}
ここまで見てきたように，IRTはさまざまな利点があります。しかし欠点がないわけではありません。

ここまでの話はすべて，項目母数が既にわかっていれば，という前提付きで進めてきました。では項目母数はどのようにして定めるのでしょう？ これはもちろん得られたデータから算出できるのですが，そのためには事前に多くの被験者から回答を集め，項目母数の値はほぼ間違いなくこれぐらいだろう，と言えるほど安定したものである必要があります。テストの場合，回答が0/1というバイナリデータで得られますから，そもそもそれほど情報がある反応ではありません。バイナリデータからその項目の特徴を安定して推定するためには，かなり多くの被験者(数千から数万単位)を集めて項目に回答させておく必要があります。テスト項目は一度使ってみないと，項目母数がどうなるかわからないというのもポイントです。

また，CATなど項目をダイナミックに組み合わせるためには，選べるぐらいさまざまな項目を準備しておかなければなりません。項目を集めたものを\keyterm{項目プール}{こうもくぷーる}{Pool of Items}といいますが，これも数千から数万の単位で用意しておく必要があります。なぜなら，テスト項目は事前に1回は使ってあるわけですから，数えるほどしか項目がなければ受験生が正答を事前に丸暗記できてしまうからです。もちろん項目プールが数千から数万あっても一度どこかで使われていますから，過去問をすべて完全に丸暗記すればその人は満点が取れてしまいます。もっともそれだけのものを覚えられるのは，ある意味学力が高いといっても差し支えないと思いますが。

日本でおこなわれる大学入学共通試験をはじめ，試験問題というのは極めて厳重な管理下に置かれ，事前にその情報が漏れることは公平性の観点から言って不適切であるとされています。しかしIRTで分析するためには，事前に項目の特徴を知っていなければならないのです。CATをつかって入学試験などができれば，カンニング対策にもなりますし，受験生は何度でもチャレンジできるので利点も多いのですが，「公平性のために新しいテストでなければならない」となるとなかなか実用化できないところがあるというのも事実です\footnote{令和2年度に大学入試センター試験から大学入学共通試験に変わりましたが，改革前の計画ではCAT化することが盛り込まれていました。しかし実際には，受験生のためのコンピュータやタブレットを準備したり，安定した通信網が必要であったり，というハード的な問題もあって見送られてしまいました。}。

ところで，心理学の場合は学力テストのように正答・誤答があるものではなく，「当てはまる」から「当てはまらない」といった軸上で多段階の反応を求めることが一般的です。テスト理論を多段階のモデルに応用できるのかと言えば，幸いにしてその答えはYESです。
\section{段階反応モデル}
リッカート法などで作られる尺度は，一般に5，7段階のものが多くあります。少ないものでは3件法\footnote{たとえばYG性格検査は3段階です。}であったり，ものによっては4，6件法であることもあります\footnote{偶数の段階にすることで，必ずどちらかの極に寄るようにして集計できます。日本人は「どちらでもない」に回答しがちな中点集中傾向があるとも言われているので，わざと肯定・否定のどちらかに寄せようという考え方です。}。しかしこれらの段階はいずれも順序尺度水準の情報しか持っておらず，そのままでは尺度値として使うことができません。シグマ法などで数値化すれば良いのですが，その手間を省いて分析する悪い習慣もあることは既に指摘した通りです。

IRTの多段階版はこうした問題に対応できる方法です。IRTの多段階モデルは大きくわけて2つあり，1つは\keyterm{段階反応モデル}{だんかいはんのうもでる}{Graded Response Model; GRM}\citep{samejima1997graded}，もうひとつは\keyterm{多段採点モデル}{ただんさいてんもでる}{Partial Credit Model}\citep{muraki1992generalized}と呼ばれています。どちらも発想は似たようなところがあり，ここではGRMについて解説します。興味がある人は，\citet{Toyoda20120824}や\citet{Katoken2014}など専門書を参考にしてください。

GRMの考え方の基本は，段階反応の背後には正規分布する連続的な潜在特性$\theta$がある，と仮定するところです。心理的な能力，学力，性質などは連続的なのですが，それが表に出てくる時は離散的(順序的)だというわけです。図\ref{fig::05_03}に図示されているように，正規分布がある\keyterm{閾値}{いきち}{threshold}(これを$b_k$と表しますが)を超えると出現する時は次のカテゴリになる，ということを考えます。図は三段階の例ですが，図から明らかなように$k$段階であれば閾値の数は$k-1$個あることになります。横軸$\theta$は心理的な態度や性質の強さだと思ってください。さてそうすると，$\theta=2.0$ぐらいであれば，ほぼ間違いなく「当てはまる」に回答することになりますし，$\theta=-2$であれば「当てはまらない」に回答するようになるはずです。このように$\theta$が大きくなればなるほど最後のカテゴリに反応する確率は上がっていきますから，ここは2PLモデルの時のようにロジスティック曲線で「当てはまる」に回答する確率を表現できるでしょう。問題は，それより下の段階に反応する確率をどのように表現するか，です。

ここである段階に反応する確率を考えるために，少し表現を改めます。すなわち，個人$i$の項目$j$に対する反応が，カテゴリ$k$\textbf{以上}になる確率として，$P^{+}_{jk}(\theta) = P(x_{ij} \ge k|\theta)$をまず考えます。ここで$k=0,1,2,...,K$とします。先ほど示したように，$k=K$，すなわち一番上のカテゴリ(ここでは「当てはまる」)に回答する確率は，2PLロジスティック関数と同じ形をしていますから，次のように表現できます。

\[ P_{jK}(\theta) = P^{+}_{jK}(\theta) = \frac{1}{1+exp(-a_j(\theta-b_{jK}))} \]

% ここで右辺の$a_j$は\keytermL{識別力}{しきべつりょく}，$b_{jK}$はカテゴリ$K$の\keytermL{困難度}{こんなんど}を表しています。
Here, the right side $a_j$ represents the \keytermL{discrimination power}, and $b_{jK}$ represents the \keytermL{difficulty level} of category $K$.
% 左辺の$P_{jK}(\theta)$は$\theta$の人が項目$j$のカテゴリ$K$に反応する確率で，それは$k$以上に反応する確率$P^{+}_{jk}(\theta)$と一致していることを表しています。
The left-hand side $P_{jK}(\theta)$ represents the probability that a person with $\theta$ will respond to category $K$ of item $j$, which indicates that it matches the probability of responding to $k$ or more $P^{+}_{jk}(\theta)$.

% 次に，もっとも低い段階に回答する確率を考えましょう。これは$\theta$が大きくなればなるほど減っていくはずで，いわばロジスティック曲線の逆のような形をするはずです。反応確率は最大でも1.0ですから，逆ということは1.0から引いてやればよいでしょう。
Next, let's consider the probability of answering at the lowest stage. This should decrease the larger $\theta$ becomes, in other words, it should take a form like the inverse of the logistic curve. Since the response probability is at most 1.0, to take the inverse, we should subtract it from 1.0.

\[ P^{+}_{j0}(\theta) = 1.0 -   \frac{1}{1+exp(-a_j(\theta-b_{j0}))} \]

問題は「どちらでもない」に反応する確率です。これは引き算で考えることができます。すなわち「どちらでもない」以上に反応する確率から，「当てはまる」以上に反応する確率を引いてやれば良いのです。

\[P_{jk}(\theta) = P^{+}_{jk}(\theta)-P^{+}_{j,k+1}(\theta) \]

% ここにあるように，段階数が増えたとしても$k$番目のカテゴリ以上に反応する確率から，$k+1$番目に反応する確率を引いてやれば，$k$番目のカテゴリに反応する確率が得られます。この計算をして描かれる曲線のことを\keyterm{項目反応カテゴリ特性曲線}{こうもくはんのうかてごりとくせいきょくせん}{Item Response Category Charactaristic Curve; IRCCC}，あるいは単に\keyterm{カテゴリ確率曲線}{かてごりかくりつきょくせん}{Category Probability Curve}と呼ばれます。IRCCCは図\ref{fig::05_03}の下段に示されています。
As shown here, even if the number of stages increases, by subtracting the probability of responding to the $k+1$th category from the probability of responding to the $k$th category or above, the probability of responding to the $k$th category is obtained. The curve drawn by doing this calculation is called the \keyterm{Item Response Category Characteristic Curve}{IRCCC} or simply the \keyterm{Category Probability Curve}. The IRCCC is shown in the lower part of Figure \ref{fig::05_03}.

\begin{figure}[ht]
    \begin{center}
        \includegraphics[width=0.8\linewidth,keepaspectratio,clip]{../images/05_GRM/Rplot05_03.png}
        \caption{正規分布と閾値(上図)とIRCCC(下図)}
        \label{fig::05_03}
    \end{center}
\end{figure}


% IRCCCを，$\theta$の値がマイナスからプラスの方向に動かしながら見ていってください。最初は当然「当てはまらない」に反応する確率が一番高いのですが，それが徐々に下がっていきます。「どちらでもない」の反応確率は徐々に増えていき，閾値$b_{j1}$で「当てはまらない」と逆転しピークを迎えることになります。その頃「当てはまる」も徐々増えていき，閾値$b_{j2}$でピークが逆転する，というようになります。ピークは逆転されても，他の確率が0になっているわけではなく，可能性は残っています。
Please observe the IRCCC as the value of $\theta$ moves from negative to positive. At first, naturally, the probability of "not applicable" is the highest, but it gradually decreases. The probability of "neither" response gradually increases and reverses with "not applicable" at the threshold $b_{j1}$, reaching a peak. Around that time, "applicable" increases gradually and the peak reverses at the threshold $b_{j2}$. Even if the peak is reversed, it does not mean that the other probabilities are zero, and the possibility remains.
% またIRCCCもICC同様に変換して，情報曲線に帰ることができます。すなわちどのあたりで鋭敏に情報を検出できるかを表現する項目反応カテゴリ情報曲線を描くこともできます。このようにして，段階反応でもその項目の特徴をデザインできるのです。
Similarly to the ICC, the IRCCC can also be converted to return to the information curve. This means that you can draw an Item Response Category Information Curve that represents where you can sensitively detect information. In this way, you can design the characteristics of that item even in the step response.

% \subsection{適切な反応段階を考える}
\subsection{Considering the Appropriate Response Stage}
% 実際の調査研究をすると，図\ref{fig::05_04}の上の段のようなIRCCCが描かれることがあります。何かおかしいところがありませんか？
When conducting actual research, an IRCCC like the one shown in the upper part of Figure \ref{fig::05_04} may be drawn. Is there anything wrong?
% これは5段階の反応モデルですが，4番目の反応カテゴリがずいぶん低く，そのピークが3番目と5番目の反応カテゴリに潰されてしまっていますね。つまり，4番目の反応がもっとも際立つシーンがないということです。言い換えるならば，これは尺度作成側が5段階だと思っていたにもかかわらず，回答者はどういう時に「やや当てはまる」と答えるのかがはっきりせず，実質4段階でしか反応していないことを表しています。
This is a five-level response model, but the fourth response category is quite low, and its peak is being squashed by the third and fifth response categories. In other words, there is no scene where the fourth reaction stands out the most. In other words, even though the scale creator believed that there were five levels, it signifies that it's unclear when the respondents would answer "somewhat applicable", and in reality, they were only responding at four levels.

% このようなIRCCCが描かれてしまう場合は，$k=3$の反応と$k=4$の反応を合わせてひとつにしてしまうなど，段階の修正を考えると良いでしょう。具体的にはデータで4と入力していたものを，3に置き換えたりします\footnote{4を5に書き換えても構いません。ヒストグラムを見て，より正規分布っぽくなるようにすると良いでしょう}。修正したのが下の図になります。このように修正しても，情報関数は変わりません。同じデータから得られる情報は同じだからです。
If such an IRCCC is drawn, it may be good to consider modifying the stages, such as combining the reaction with $k=3$ and the reaction with $k=4$ into one. Specifically, you replace the ones that were input as 4 in the data with 3\footnote{You can also rewrite 4 to 5. It is a good idea to look at the histogram and make it look more like a normal distribution}. The figure after the modification is shown below. Even if you make such adjustments, the information function does not change. This is because the information obtained from the same data is the same.

% このように5段階，7段階を設定して回答者に無理やり回答を求めても，分析するとカテゴリのピークが潰れていることがあります。回答者の反応しやすいカテゴリ数を丁寧に設計してやることが重要です。もちろん無分別に尺度値をつけて，そのまま分析するのはもっとも不適切な方法です。
Even if you force respondents to respond by setting five or seven stages, there may be instances where the peak of the categories is crushed upon analysis. It is important to carefully design the number of categories that respondents can easily react to. Of course, recklessly assigning scale values and analyzing them as they are is the most inappropriate method.
\begin{figure}[ht]
	\begin{center}
		\includegraphics[width=0.8\linewidth,keepaspectratio,clip]{../images/05_GRM/Rplot05_04.png}
		\caption{適切な反応段階をデザインする}
		\label{fig::05_04}
	\end{center}
\end{figure}




% \section{因子分析の歴史と展開}
\section{History and Development of Factor Analysis}

% ところで，因子分析モデルもテスト理論も，潜在変数モデルとしては同じでいずれも古典的テスト理論からの発展系なのでした。
By the way, both factor analysis models and test theory are the same as latent variable models, and both are developments from classical test theory.
% 因子分析モデルは多段階反応が一般的で，多因子モデルで「潜在的な(心理学的な)構造はどうか」ということを問題にします。ここでの目的は全体に共通する要素やその構造であり，何種類に別れて要素間の関係はどうなっているのか、というところが中心的関心事になります。
Factor analysis models typically involve multi-step responses, and multi-factor models deal with the issue of "what is the latent (psychological) structure?". The concern here is about common elements and their structure overall: the central interest lies in how many kinds of relationships exist between elements.
% 一方，項目反応理論はバイナリ反応が一般的で，因子数はひとつです。学力テストはその学力が測定できていることが重要で，因子の構造よりも因子得点をより精緻に推定できることの方が重要だからです。因子分析の言葉で言えば，因子得点をより精緻に表現しようとしているわけです。
Meanwhile, item response theory commonly involves binary responses and a single factor. It is crucial that ability tests are able to measure that ability, which is more important than the structure of the factors, because what matters most is that it can estimate factor scores more precisely. In the language of factor analysis, it's about expressing factor scores more intricately.

% さて，GRMは，項目反応理論の多段階モデルでした。実はGRMは因子分析モデルの発展系でもあるのです。
Now, the GRM was a multistage model of the item response theory. In fact, the GRM is also a developed version of the factor analysis model.
% 因子分析は相関係数のモデルであったことを思い出してください。因子分析モデル自体は$z_{ij}$から始まっていましたが，変数同士の共分散$r_{ij}$を考えるといくつかの仮定から\keytermL{因子負荷量}{いんしふかりょう}だけのモデルに簡略化され，推定できるようになるのでした\footnote{具体的にどうやって因子負荷量を算出するかは次回以降のお楽しみです。}。この標準化された共分散，すなわち相関係数はピアソンの積率相関係数とも言われ，間隔尺度水準以上の数値を使って計算されます。多段階の反応は順序尺度水準ですから，相関係数を計算するのは不適切で，そのまま因子分析することはできません\footnote{できないのですが，尺度値に変換することもなく機械的に分析してしまう悪い習慣が蔓延っているのは何度も指摘している通りです。くどいと思われるかもしれませんが，私は憤っているのです。}。では\keytermL{順序尺度水準}{じゅんじょしゃくどすいじゅん}の相関係数がないのかといわれると，あります。
Remember that factor analysis is a model of correlation coefficients. The factor analysis model itself started from $z_{ij}$, but when you consider the covariance $r_{ij}$ between variables, it is simplified to a model of only \keytermL{factor loadings} due to several assumptions, and it becomes estimable\footnote{How to calculate the factor loadings will be fun for next time.}. This standardized covariance, also known as Pearson's product-moment correlation coefficient, is calculated using values at the interval scale level or above. Multi-step reactions are at the ordinal scale level, so it is inappropriate to calculate correlation coefficients and it is not possible to perform factor analysis directly\footnote{Although it cannot be done, I have pointed out many times that the bad habit of mechanically analyzing without converting to scale values is rampant. Although you may get tired of hearing this, I am indignant.}. Then, if asked if there is no correlation coefficient at the \keytermL{ordinal scale level}, it exists.

% 順序尺度水準の変数$\times$順序尺度水準の変数の相関は\keyterm{ポリコリック相関係数}{ぽりこりっくそうかんけいすう}{polychoric correlation}といいます。
The correlation of an ordinal scale level variable $\times$ ordinal scale level variable is called the polychoric correlation coefficient.
% 順序尺度水準の変数$\times$間隔尺度水準の変数の相関は\keyterm{ポリシリアル相関係数}{ぽりしりあるそうかんけいすう}{polyserial correlation}といいます。
The correlation between variables at ordinal scale level and interval scale level is called the polyserial correlation coefficient.
% ついでにバイナリ変数$\times$バイナリ変数の相関係数は\keyterm{テトラコリック相関係数}{てとらこりっくそうかんけいすう}{tetracholic correlation}といいます。
By the way, the correlation coefficient of the binary variable x binary variable is called the "tetracholic correlation".

% これらの相関係数はいずれも，順序(あるいはバイナリ)変数の背後に連続体があると考え，潜在的な連続体$\times$潜在的な連続体の相関係数を連続体のカテゴリが変わる閾値とともに推定するのです(図\ref{fig::05_05})。
All of these correlation coefficients assume that there is a continuum behind the ordinal (or binary) variables, and estimate the correlation coefficient of the latent continuum × latent continuum as the threshold changes for categories of the continuum (Figure \ref{fig::05_05}).
\begin{figure}[ht]
	\begin{center}
		\includegraphics[width=0.8\linewidth,keepaspectratio,clip]{../images/05_GRM/Slide05_01.png}
		\caption{ポリコリック相関係数のイメージ}
		\label{fig::05_05}
	\end{center}
\end{figure}


% こうして推定された相関係数を使って因子分析を行うと，順序尺度水準に適した因子分析を行うことができます。この方法をとくに\keyterm{カテゴリカル因子分析}{かてごりかるいんしぶんせき}{categorical factor analysis}といいます。もっともこの名前を覚えておく必要はありません。\textbf{カテゴリカル因子分析は段階反応モデルと数学的に等価である}ことがわかっています。GRMをすることはカテゴリカル因子分析をしていることと同じ，なのです。
By performing a factor analysis using the estimated correlation coefficients, we can conduct a factor analysis suitable for ordinal scale levels. This method is particularly referred to as a \keyterm{Categorical Factor Analysis}. However, it is not necessary to remember this name. It is known that \textbf{Categorical Factor Analysis is mathematically equivalent to the Graded Response Model}. So, conducting a GRM is the same as conducting a Categorical Factor Analysis.

% もっともGRMは項目反応理論の系列ですから，単因子構造を仮定しています。複数の因子を想定するカテゴリカル因子分析に対応するのは，正確には\keyterm{多次元項目反応理論}{たじげんこうもくはんのうりろん}{Multidimensional Item Response Theory}といいます。しかし数学的・技術的には同じであり，カテゴリカル因子分析をした結果からIRCCCを描くこともできますし，実際に分析するソフト上では使用する変数がカテゴリカル(順序尺度水準)であることを指定するだけです。われわれユーザはもはや悩む必要はなく，ただただデータに適した分析をするだけで良いのです。
Sandly, the GRM is a series of item response theories, so it assumes a single-factor structure. Strictly speaking, the corresponding categorical factor analysis, which assumes multiple factors, is called the Multidimensional Item Response Theory. However, mathematically and technically they are the same, and it is possible to draw an IRCCC from the results of categorical factor analysis, and in actual analysis software, all that needs to be specified is that the variables used are categorical (ordinal scale level). We users no longer need to worry, we just need to do an analysis appropriate for the data.

% \subsection{系譜の違いはどこに関係するか}
\subsection{Where do Differences in Lineage Matter?}

% 因子分析モデルと項目反応理論が，カテゴリカル因子分析として概念的に統合されることを話してきました。本質的にはこのように違いがないのですが，系譜の違い，出自の違いはそれぞれの利用される文脈で，何を強調するかに影響してくることがあります。
We have been discussing how factor analysis models and item response theory are conceptually integrated as categorical factor analysis. Essentially, there is no difference, but differences in lineage and origin may impact what is emphasized in the context in which each is used.

% たとえば因子分析の文脈では，共通性が低い項目は削除し，綺麗な因子構造を目指そうという考え方があります。尺度作成の中で1つの項目は1つの因子に対応しているべきであるという考え方があり(\keyterm{単純構造の原理}{たんじゅんこうぞうのげんり}{Principle of Simple Structure}といいます)，もし1つの項目が複数の因子の影響を受けているようであれば，「美しくないので」削除されることがあります。因子分析は知能，性格，態度の研究で展開されてきたため，美しい「構造」を見つけ出すことに狙いがありますから，この美しさを損なうもの(項目)は取り除く，という方向に行きがちです。
For example, in the context of factor analysis, there is a concept of removing items with low commonality and aiming for a neat factor structure. There is also a concept that each item should correspond to one factor in scale creation (this is called the Principle of Simple Structure), and if one item appears to be influenced by multiple factors, it may be removed because it's not 'beautiful'. Factor analysis has been used in research on intelligence, personality, and attitudes, aiming to find a beautiful 'structure', so it tends to go in the direction of removing things (items) that spoil this beauty.

% 一方で，項目反応理論はテスト理論の生まれです。もちろん回答者の能力や特性を測定するのに優れた項目とそうでない項目，という峻別はしますが，中でも「この項目は測定能力の偏差値30程度の回答者を測定するのに適している」とか，「偏差値70程度の回答者を測定するのに適している」と判断できます。偏差値30や70を測定するのに適した項目とは，非常に簡単(ほとんどの人が正答する)だったり，非常に難易度が高い(ほとんどの人が誤答する)項目です。心理尺度でいうと床効果，天井効果がみられる項目とされる，どちらかに偏った分布をもつ項目です。しかし，それを捨てるということにはならず，どのような項目でも回答者の能力を推定するための情報量がゼロではない，という考えから，さまざまな項目をどんどんためていく方向にいきがちです。項目反応理論の文脈では，あらゆる人に対する測定を準備しておく必要があり，むしろ回答者の特性にあわせて設問の方を選んだり，事前の項目特性から，前もってテストの項目構成をデザインする，ということを目的とするのです。
On the other hand, the Item Response Theory comes from the theory of testing. Of course, there is a clear distinction between excellent items and not-so-excellent items in measuring the abilities and characteristics of the respondents, but among them, you can judge that "this item is suitable for measuring a respondent with a deviation value of about 30" or "it's suitable for measuring a respondent with a deviation value of about 70". Items that are appropriate for measuring deviation values of 30 or 70 are either extremely easy (that most people answer correctly) or extremely difficult (that most people answer incorrectly). In the language of psychological scales, these are items with skewed distributions which tend to show floor or ceiling effects. However, the thought is instead of discarding them, no kind of item contains zero information for estimating the respondent's ability, so there's a tendency to pile up various items. In the context of Item Response Theory, it's necessary to prepare measurements for all people, and instead the aim is to select questions according to the characteristics of the respondent, or to design the test item composition in advance from the item characteristics beforehand.

% このように使われるシーンによって，「構造」か「機能(=得点)」のどちらに注目するかが変わり，結果的に実践の方針がちがってくることもあるのです。図\ref{fig::05_06}にこの心理学的系譜(左ルート)とテスト理論的系譜(右ルート)の流れを描いてみました。
The focus may change between "structure" and "function (= scoring)" depending on the scene it is used, and as a result, the policy of practice may also change. I've drawn the flow of this psychological lineage (left route) and test theory lineage (right route) in Figure \ref{fig::05_06}.

% ポイントは「最終的には同じところに辿り着く」という点ですから，歴史的流れや個々のモデルの細かい数式を完全に理解していなくてもいいかもしれません。それよりは，心理学者として，あるいはテストをする側として，回答者に無理のない，それでいて必要な程度に精緻な情報が得られる適切な分析方法を選択できるようになることが重要です。
The point is that "we will eventually arrive at the same place," so you may not need to fully understand the historical flow and detailed equations of each model. Rather, as a psychologist, or as a test administrator, it's important to be able to choose an appropriate analysis method that provides detailed information to the extent necessary, without overburdening the respondent.
\begin{figure}[ht]
	\begin{tikzpicture}
		\node[draw,align=center] (a) at (5,0)
		{古典的テスト理論\\$X=t+e$};
		\node(P0) at (5,0.9)[above] {見えないものを測定・数値化したい};
		\node[draw,align=center] (b) at (0,-2.5) {探索的因子分析\\$z_{ij}=\sum a_{jm}f_{im} + d_jU_{ij}$};
		\node[draw,align=center] (c) at (10,-2.5) {項目反応理論\\ $p_j(\theta) = \frac{1}{1+exp(-1.7a_j(\theta_i -b_j))}$};
		\node[draw] (d) at (5,-5) {カテゴリカル因子分析};
		\draw[->] (a) --(b);
		\draw[->] (a) --(c);
		\draw[->] (b) --(d);
		\draw[->] (c) --(d);
		\node (P1) at (1.3,-1.25)[above] {構造に興味がある};
		\node (P2) at (8.7,-1.25)[above] {機能に興味がある};
		\node (P3) at (2.5,-4.2)[left,align=center] {順序尺度水準へ\\得点の精緻化};
		\node (P4) at (7.5,-4.2)[right,align=center] {多段階反応へ\\多次元に拡張};
	\end{tikzpicture}
	\caption{理論・モデルの流れ。左側のルートが心理学的系譜，右側のルートがテスト理論的系譜}
	\label{fig::05_06}
\end{figure}


% \section{課題}
\section{Task}
% \paragraph{信頼性についての考え方} 古典的テスト理論，因子分析モデル，項目反応理論，それぞれの信頼性の考え方を数式とともに自分のことばで説明できるようになろう。
\paragraph{Understanding Reliability} Let's try to explain the concepts of reliability in classical test theory, factor analysis model, and item response theory using mathematical formulas in our own words.
% \paragraph{段階反応モデルのIRCCC} 多段階モデルのIRCCCを，関数を描画するソフトなどを使って自分で描いてみよう。
\paragraph{IRCCC of Stepwise Reaction Model} Try drawing the IRCCC of the multi-step model yourself, using software for plotting functions.
\end{document}
